\chapter{Arquitetura do Sistema}

A rede possui arquitetura hierárquica. Clusters locais, com linhas de base de dezenas de quilômetros, atendem iluminadores regionais e cenários de baixa altitude. Clusters regionais ampliam a cobertura, enquanto linhas de base longas podem apoiar observações em HF, alvos de alta altitude e objetos espaciais quando há visada comum e iluminador compatível.

Cada estação contém plataforma radome, cadeias RF multifaixa, processamento local, sincronismo e calibração, transporte óptico, condicionamento de energia, gestão térmica e interface de supervisão. A fusão central associa detecções por forma de onda, tempo, frequência, Doppler, polarização e consistência geométrica. Deve distinguir observações diretas de emissores, reflexões biestáticas, interferência e hipóteses de rastreamento.

\begin{figure}[H]
\centering
\includegraphics[width=0.95\textwidth]{pt-BR/fig01_arquitetura_rede.png}
\caption{Arquitetura distribuída de sensoriamento: o polígono superior representa um iluminador ou emissor, os três domos são nós receptores independentes, os caminhos verdes transportam observações calibradas e o bloco roxo inferior representa a fusão central.}
\label{fig:fig01_arquitetura_rede}
\end{figure}
